% ============================================================
%  SECTION 1: INTRODUCTION
% ============================================================

\section{Introduction}

Honey is one of the oldest food products in human history, and at
the same time one of the most frequently adulterated commodities
on~the modern food market. According to data from the European Food
Fraud Database (FoodFraudNetwork), honey ranks third on the list of
products with the highest adulteration risk, behind only milk
and~olive oil~\cite{apimondia_fraud_v2}. In~2023, the total volume
of~honey imports into the European Union exceeded 360,000~tonnes
at~a~value close to EUR~1~billion~\cite{cbi2024}, while the EU's own
production amounted to approximately 286,000~tonnes, covering
only 60\%~of domestic demand~\cite{pollinatorhub2025}.
The EU's structural production deficit makes the market largely
dependent on supplies from third countries, including China (35\%
of~imports), Ukraine (31\%), and Argentina (12\%).
The scale of this deficit, combined with the significant price
disparity between authentic honey and syrup substitutes, creates
economic incentives for fraudulent practices~\cite{apimondia_fraud_v2}.

In the face of these challenges, Apimondia (the International
Federation of Beekeepers' Associations) plays a key role. This
organization, uniting more than 80 national organizations and
headquartered in Rome, constitutes a global forum for cooperation
aimed at promoting scientific, technical, and economic progress
in beekeeping~\cite{apimondia_web}.

The Apimondia position published in~2025 by the Working Group
on~Adulteration of Bee Products constitutes an important voice
in~the debate on product integrity~\cite{garcia2025}.
The author identifies the systemic production of immature honey,
defined as the deliberate collection of honey before the completion
of the biological ripening process in the hive, combined with its
industrial dehydration, as a separate and serious form of adulteration,
distinct from classical practices of adulteration with sugar
syrups~\cite{garcia2025}.
The significance of this position is considerable: Apimondia, as the
world's largest federation of beekeeping organizations, has real
influence on the shaping of both international standards and control
practices applied by sanitary inspectorates in EU member states.

This article does not question the existence or seriousness of
systemic immature honey production as a phenomenon violating the
principles of fair food trade. \textbf{The aim of this paper is,
however, to indicate that} the Apimondia position, perhaps
inadvertently, makes an excessive generalization, equating three
qualitatively distinct phenomena:
(1)~systemic, deliberate production of immature honey for economic
gain; (2)~early honey harvesting under conditions forced by climatic
or technological factors; and (3)~the presence of osmophilic yeasts
as natural components of honey~\cite{garcia2025,munitis1976}.
The consequence of this generalization is the creation of a
disqualifying criterion based on the presence of yeasts, which
remains in contradiction with both applicable regulatory standards
and the current state of microbiological
knowledge~\cite{codex_stan12, rodriguezmachado2024}.

Analysis of applicable international standards, including Codex
Alimentarius CODEX STAN 12-1981, EU Directive 2001/110/EC, and the
draft ISO/DIS 24607:2025 standard, leads to the conclusion that
none of these documents establishes a limit on yeast cell count
(CFU/g) as a criterion for assessing the quality or authenticity of
honey~\cite{codex_stan12}.
The only legal microbiological criterion is \emph{fermentation as an
active state}, not merely the presence of microorganisms capable of
fermentation~\cite{codex_stan12}.
Meanwhile, the fermentative activity of osmophilic yeasts, including
the species \textit{Zygosaccharomyces rouxii} and \textit{Z.~richteri},
endemically present in honey of all botanical and geographical types,
is biologically impossible at water activity $a_w < 0{,}60$,
corresponding to honey moisture content of $\leq 20\%$~\cite{zamora2006}.

The above premises lead to the three-part thesis of this article:
first, the presence of osmophilic yeasts in honey meeting the moisture
criterion of $\leq 20\%$ is normatively neutral, since water activity
$a_w < 0{,}60$ prevents the growth and fermentation of these
microorganisms; second, no applicable national or international
standard penalizes the mere presence of yeasts without active
fermentation, and therefore the determination of their count by the
plate method (CFU/g) alone constitutes neither a legal nor a scientific
basis for disqualifying the product from trade; third, the concept of
\emph{functional maturity}, understood as integrated microbiological,
chemical, and biological stability, constitutes a more precise and
operationalizable tool for honey quality assessment than the sole
morphological criterion of comb cell
capping~\cite{zhang2021,guo2020,sun2021}.

The article is a review and covers the following areas:
(i)~analysis of international regulatory frameworks concerning
the microbiological parameters of honey;
(ii)~biology and metabolomics of the honey ripening process
as a continuum;
(iii)~microbiology of osmophilic yeasts in the context of water
activity and fermentation risk;
(iv)~diversity of global beekeeping production systems and their
implications for standardization;
(v)~analysis of the EU honey market and the regulatory consequences
of overly strict interpretations; and
(vi)~a proposal for a functional maturity model with recommendations
for standardization bodies and Apimondia.
